\section{Experiments}
\label{sec:experiments}

\subsection{Experimental Setup}
\label{sec:setup}

\paragraph{Datasets.}
We evaluate on three English IE benchmarks spanning a difficulty gradient.
\textbf{SciERC} \cite{luan2018scierc} contains 551 test instances from scientific abstracts annotated for joint NER (6 entity types) and RE (7 relation types), representing a complex schema with overlapping structures.
\textbf{CoNLL-2003} \cite{tjongkimsang2003conll} provides 3{,}453 test sentences from newswire with 4 entity types, serving as a high-accuracy benchmark where models achieve ${>}$91\% micro F1.
\textbf{WNUT-17} \cite{derczynski2017wnut} consists of 1{,}287 test instances from social media with emerging entities, used to probe method boundaries where model capability is insufficient.
This gradient---from near-solved (CoNLL) through challenging (SciERC) to infeasible (WNUT-17)---enables systematic characterization of how signal effectiveness varies with task difficulty (\S\ref{sec:complexity_robustness}).

\paragraph{Models.}
Our primary model is \textbf{Qwen3-8B}, fine-tuned with LoRA ($r{=}32$) separately per dataset.
For cross-model validation, we evaluate \textbf{LLaMA-3.1-8B-Instruct} to assess whether structural confidence signals transfer across architectures and training regimes.

\paragraph{Sampling and Evaluation.}
For each test instance, we draw $N$ independent samples at temperature $T{=}1.0$, with $T{=}0.7$ for robustness analysis.
We use $N{=}8$ for both SciERC and CoNLL in the primary evaluation, and $N{=}16$ for SciERC scaling experiments.
Quality estimation is measured along two axes: \textit{rank correlation} with instance-level F1 via Spearman $\rho$ and Kendall $\tau$, and \textit{binary discrimination} via AUROC (instances above vs.\ below median F1).
Downstream utility is assessed through \textbf{best-of-$N$ selection F1}: the F1 obtained when each signal selects one sample per instance.
Instance-level F1 uses exact span matching for NER and head-tail-label matching for RE, where a relation is correct only when both argument entities and the relation label match the gold annotation.
The greedy decoding baseline ($T{=}0$, $N{=}1$) establishes performance without test-time scaling.

\paragraph{Confidence Signals.}
We compare five signals from three families (detailed in \S\ref{sec:signals}).
\textit{Structural} signals operate on extracted entity/relation sets: \textbf{Voting Confidence} aggregates per-extraction frequency across samples, and \textbf{Soft Jaccard (SJ)} measures pairwise set similarity \cite{jaccard1901distribution}.
\textit{Surface} signals operate on string representations: \textbf{Fleiss' $\kappa$ (FK)} quantifies inter-sample agreement treating samples as independent annotators \cite{fleiss1971agreement}, and \textbf{Exact Match (EM)} computes the fraction of identical sample pairs.
The \textit{model-internal} signal is \textbf{log-probability (logprob)}, the mean token-level log-likelihood.

\paragraph{Implementation.}
Inference uses vLLM with XGrammar constrained decoding \cite{willard2023guided} to enforce valid JSON output structure.
A span realignment post-processing step corrects character-level offset mismatches between generated and gold annotation spans.
All confidence signals are computed from the $N$ samples without access to gold labels.

%% ============================================================
\subsection{Quality Estimation Results}
\label{sec:qe_results}

\begin{table*}[t]
\centering
\small
\caption{Quality estimation across tasks and models (Spearman $\rho$ and AUROC). All results use $N{=}8$ samples at $T{=}1.0$. Best value per column in \textbf{bold}. SciERC evaluates joint NER (6 types) and RE (7 types); CoNLL evaluates NER (4 types) with two model families. RE results are preliminary (single dataset).}
\label{tab:quality_estimation}
\begin{tabular}{ll cc cc cc cc}
\toprule
& & \multicolumn{2}{c}{\textbf{SciERC NER}} & \multicolumn{2}{c}{\textbf{SciERC RE}} & \multicolumn{2}{c}{\textbf{CoNLL (Qwen)}} & \multicolumn{2}{c}{\textbf{CoNLL (LLaMA)}} \\
\cmidrule(lr){3-4} \cmidrule(lr){5-6} \cmidrule(lr){7-8} \cmidrule(lr){9-10}
Signal & Family & $\rho$ & AUC & $\rho$ & AUC & $\rho$ & AUC & $\rho$ & AUC \\
\midrule
Voting Conf. & Structural & \textbf{.379} & .684 & \textbf{.350} & .694 & .424 & .789 & .426 & .709 \\
Soft Jaccard  & Structural & .360 & .692 & .250 & .670 & .436 & .801 & \textbf{.431} & .704 \\
Fleiss' $\kappa$ & Surface & .267 & .629 & .275 & .609 & .400 & .774 & .428 & .703 \\
Exact Match  & Surface    & .295 & \textbf{.749} & .134 & \textbf{.706} & \textbf{.451} & \textbf{.816} & \textbf{.431} & .705 \\
Log-prob     & Model      & .205 & .578 & .266 & .601 & .316 & .753 & .307 & \textbf{.761} \\
\bottomrule
\end{tabular}
\end{table*}

Table~\ref{tab:quality_estimation} presents rank correlation ($\rho$) and AUROC for five confidence signals across tasks and models.

\paragraph{Structural signals dominate rank correlation on complex schemas.}
On SciERC NER, structural signals lead rank correlation: Voting Confidence achieves $\rho{=}0.379$ with SJ close behind ($\rho{=}0.360$), both substantially outperforming surface and model-internal alternatives ($\rho{\leq}0.295$).
Voting's edge over SJ reflects its extraction-level granularity: it rewards entities that appear consistently while penalizing hallucinated or missed ones, whereas SJ compares entire sets and is sensitive to both overlap and total extraction count.
The structural lead persists on RE, where Voting ($\rho{=}0.350$) again leads.
However, AUROC rankings diverge: Exact Match achieves the highest AUROC on both NER (0.749) and RE (0.706) despite its lower $\rho$, revealing a ranking-discrimination dissociation we analyze below (\S\ref{sec:complexity_robustness}).
All signals weaken on RE: 40\% of instances receive zero F1, creating a floor effect that compresses the discriminable range.
FK ($\rho{=}0.275$) exceeds SJ ($\rho{=}0.250$) on RE---a reversal from NER showing no single structural signal dominates across both tasks.

\paragraph{Signal rankings shift with task difficulty.}
On CoNLL NER, where greedy decoding achieves ${>}$91\% micro F1, the signal landscape changes substantially.
Exact Match leads with $\rho{=}0.451$ and AUROC${=}0.816$, surpassing SJ ($\rho{=}0.436$) and Voting ($\rho{=}0.424$).
When most samples are near-correct, binary agreement captures the remaining variance more efficiently than set-level similarity.
This reversal---EM${>}$SJ on easy tasks vs.\ SJ${>}$EM on hard tasks---demonstrates that optimal signal choice depends on intrinsic task difficulty.

\paragraph{Cross-model stability.}
Comparing Qwen and LLaMA on CoNLL, SJ shows remarkable stability: $\rho{=}0.436$ (Qwen) vs.\ $\rho{=}0.431$ (LLaMA), a gap of only 0.005.
This model-agnostic property is expected for structural signals that operate on outputs rather than internal model states, and provides the strongest evidence that SJ generalizes across architectures.
The SJ--FK gap also narrows with schema simplification: $\Delta\rho{=}0.092$ on SciERC NER (6 entity types) vs.\ $\Delta\rho{=}0.003$ on CoNLL with LLaMA (4 types), suggesting that set-level structure contributes less when the type inventory is small.

\paragraph{The logprob paradox.}
Log-probability reveals a dissociation between ranking and classification.
On CoNLL with LLaMA, it achieves the highest AUROC (0.761) yet the lowest $\rho$ (0.307) among all signals.
This pattern indicates that logprob assigns well-separated scores to easy vs.\ hard instances (strong binary discrimination) but provides little resolution within each group (weak ranking), making it better suited for binary filtering than continuous quality estimation.
We return to this distinction in \S\ref{sec:complexity_robustness} and visualize it in Figure~\ref{fig:signal_heatmap}.

\begin{figure}[t]
\centering
\includegraphics[width=\columnwidth]{figures/fig_signal_heatmap_1024}
\caption{Signal quality heatmap across tasks and models. Color intensity encodes Spearman $\rho$; structural signals (VC, SJ) dominate on complex schemas (SciERC) while surface signals (EM) lead on simpler tasks (CoNLL).}
\label{fig:signal_heatmap}
\end{figure}

\paragraph{Oracle headroom and practical ceilings.}
On CoNLL, oracle best-of-$N$ selection yields only +4.2~pp (Qwen) and +1.8~pp (LLaMA) over greedy.
When the model already operates near its performance ceiling, confidence signals serve as \textit{reliability assurance}---confirming that outputs are trustworthy---rather than enabling selective prediction gains.
We therefore focus selection analysis (\S\ref{sec:selection_gap}) on SciERC, where oracle headroom exceeds 16~pp.

%% ============================================================
\subsection{The Correlation-Selection Gap}
\label{sec:selection_gap}

\begin{table}[t]
\centering
\small
\caption{Best-of-$N$ selection F1 on SciERC ($N{=}8$, $T{=}1.0$). Best non-oracle result per column in \textbf{bold}. EM produces binary agreement without sample ranking. RE results are preliminary (single dataset, single model).}
\label{tab:selection}
\begin{tabular}{lcc}
\toprule
Method & NER F1 & RE F1 \\
\midrule
Oracle & 0.769 & 0.561 \\
\midrule
Log-prob           & 0.609 & \textbf{0.397} \\
Ensemble (SJ+Vot.) & 0.604 & 0.367 \\
Soft Jaccard       & 0.604 & 0.368 \\
Voting             & 0.596 & 0.338 \\
Exact Match        & ---   & ---   \\
\midrule
Greedy ($T{=}0$)   & \textbf{0.609} & 0.374 \\
Random             & 0.580 & 0.338 \\
\bottomrule
\end{tabular}
\end{table}

Despite strong quality estimation, confidence signals yield limited gains when used for best-of-$N$ selection.
Table~\ref{tab:selection} compares selection F1 on SciERC ($N{=}8$).

\paragraph{NER: no signal surpasses greedy.}
The oracle upper bound (0.769) reveals 16.0~pp of headroom above greedy (0.609).
Yet no confidence signal closes this gap.
Log-probability (0.609) merely ties greedy; the best quality estimator, Voting ($\rho{=}0.379$), selects to only 0.596, \textit{below} greedy.
An ensemble of SJ and Voting (AUROC${=}0.735$) also fails, achieving 0.604.
The greedy--logprob tie is revealing: log-probability selects the sample with highest model confidence, which on average matches the deterministic greedy output, implying that the $N$ samples cluster at a similar quality level with insufficient variance for any signal to exploit.
The Random baseline (0.580) confirms that not all samples are equal---uniform selection loses 2.9~pp vs.\ greedy---but the discriminative structure captured by confidence signals adds no value beyond greedy's deterministic choice.

\paragraph{RE: logprob uniquely succeeds.}
For RE, log-probability is the only signal to exceed greedy (0.397 vs.\ 0.374, +2.3~pp).
Structural signals trail or match greedy: SJ selects to 0.368 and Voting to 0.338.
The oracle gap of 18.7~pp (0.561 vs.\ 0.374) indicates substantially more per-instance variance than NER.
Log-probability's success here may stem from its per-token nature: unlike instance-level consistency signals, logprob evaluates each sample independently, enabling it to distinguish a particularly successful extraction from its peers even when all extractions share similar structural consistency.

\begin{figure}[t]
\centering
\includegraphics[width=\columnwidth]{figures/fig_selection_gap}
\caption{The correlation-selection gap on SciERC ($N{=}8$). Left: moderate QE correlation (Spearman $\rho$ up to 0.38). Right: negligible selection gains (best-of-$N$ F1 at or below greedy). Dashed line indicates greedy baseline.}
\label{fig:selection_gap}
\end{figure}

\paragraph{Quality estimation $\neq$ output selection.}
We attribute this gap (Figure~\ref{fig:selection_gap}) to a fundamental distinction: instance-level confidence signals predict \textit{which inputs} the model handles reliably, not \textit{which output} among $N$ candidates is best for a given input.
When per-sample variance is low relative to instance-level difficulty---as in NER, where the model often produces $N$ similar extractions---even perfect quality estimation cannot discriminate among near-identical candidates.

\paragraph{Implications for test-time scaling.}
This correlation-selection gap has direct consequences for test-time scaling in IE \cite{brown2024monkeys, snell2024scaling}.
Scaling inference compute through repeated sampling assumes a selector can identify the best sample.
Our results show that for structured prediction, existing confidence signals serve as quality \textit{estimators}---useful for flagging unreliable predictions---but not as effective \textit{selectors}.
Closing this gap likely requires signals that capture per-sample rather than per-instance variation, such as trained verifiers \cite{cobbe2021verifiers} or pairwise comparison models.

%% ============================================================
\subsection{Task Complexity and Signal Robustness}
\label{sec:complexity_robustness}

\paragraph{Task difficulty as a boundary condition.}
Signal effectiveness varies systematically with task difficulty.
On CoNLL (greedy micro F1${>}$0.91), EM leads all signals ($\rho{=}0.451$); on SciERC NER (micro F1$\,{\approx}\,$0.66), SJ overtakes EM ($\rho{=}0.357$ vs.\ 0.257); on SciERC RE, where 44\% of instances score zero F1, all signals converge to $\rho{\approx}0.25$.
The NER--RE difference is statistically confirmed: Fisher's $z$-test comparing SJ $\rho$ between NER ($\rho{=}0.357$) and RE ($\rho{=}0.211$) yields $z{=}2.31$, $p{=}0.021$, confirming statistically significant task-dependent degradation.

\paragraph{Per-relation diagnostic.}
RE signal degradation is relation-specific.
\textsc{Feature-Of} has near-zero correlation ($\rho{=}0.004$), while \textsc{Hyponym-Of} achieves $\rho{=}0.646$.
The primary bottleneck is entity span errors, which account for 55.5\% of RE false negatives---a prerequisite failure that consistency signals cannot detect, as they measure agreement among equally incorrect extractions.

\paragraph{Capability stratification.}
SJ effectiveness scales with structural complexity: on SciERC instances with 6+ entities, $\rho{=}0.547$, compared to $\rho{=}0.415$ for 1--2 entities.
Similarly, RE instances with ${\geq}3$ relation triples ($n{=}81$) yield substantially stronger correlations (Voting $\rho{=}0.630$, SJ $\rho{=}0.502$), demonstrating that within the model's capability range, increased output complexity amplifies structural signal advantage.
The SJ advantage over logprob remains stable across sentence-length quartiles ($\Delta\rho{\approx}0.08$--$0.20$), and three independent complexity proxies (entity count, sentence length, relation density) show weak mutual correlation ($|r|{<}0.15$), indicating a genuinely multi-dimensional benefit rather than a single confound.

\paragraph{Method boundary: WNUT-17.}
WNUT-17 establishes where the approach breaks down.
SJ achieves $\rho{=}0.062$ ($p{=}0.106$, not statistically significant), with 54\% of instances at zero F1.
Log-probability shows a strong \textit{negative} correlation ($\rho{=}{-}0.411$): the model is confidently wrong on novel entities, a calibration failure that no post-hoc consistency signal can remedy.
This is a boundary of the method, not a limitation of the datasets.

\begin{table}[t]
\centering
\small
\caption{Signal robustness across temperature and sample size on SciERC (Spearman $\rho$). $^\dagger$Mean over 3 seeds ($\sigma{=}0.025$; individual: 0.428, 0.424, 0.384). RE results are preliminary (single dataset).}
\label{tab:robustness}
\begin{tabular}{llcc}
\toprule
& & NER $\rho$ & RE $\rho$ \\
\midrule
\multicolumn{4}{l}{\textit{Temperature} ($N{=}8$)} \\
\quad SJ & $T{=}0.7$ & 0.388 & 0.234 \\
\quad SJ & $T{=}1.0$ & 0.357 & 0.211 \\
\quad FK & $T{=}0.7$ & 0.321 & 0.184 \\
\quad FK & $T{=}1.0$ & 0.265 & 0.253 \\
\midrule
\multicolumn{4}{l}{\textit{Sample size} ($T{=}1.0$)} \\
\quad SJ & $N{=}8$  & 0.357 & 0.211 \\
\quad SJ & $N{=}16$ & 0.412$^\dagger$ & 0.281 \\
\quad FK & $N{=}8$  & 0.265 & 0.253 \\
\quad FK & $N{=}16$ & 0.280 & 0.279 \\
\bottomrule
\end{tabular}
\end{table}

\paragraph{Robustness analysis.}
Table~\ref{tab:robustness} reports signal stability across temperature and sample size on SciERC.
\textit{Temperature.}
SJ NER $\rho$ shifts from 0.388 ($T{=}0.7$) to 0.357 ($T{=}1.0$), a difference of 0.031 whose bootstrap confidence interval includes zero.
FK shows mixed behavior: NER $\rho$ drops by 0.056, while RE $\rho$ increases by 0.069, likely reflecting reduced diversity at lower temperature interacting differently across tasks.
Across all comparisons, the SJ--FK relative ordering is preserved at both temperatures on NER (0.388 vs.\ 0.321 at $T{=}0.7$; 0.357 vs.\ 0.265 at $T{=}1.0$), indicating a stable signal hierarchy despite numerically small temperature effects.

\textit{Sample size.}
Increasing from $N{=}8$ to $N{=}16$ yields meaningful gains for SJ on NER: $\rho$ improves from 0.357 to $0.412{\pm}0.025$ (three seeds: 0.428, 0.424, 0.384).
FK benefits less ($\rho{\approx}0.280$ at $N{=}16$ vs.\ 0.265 at $N{=}8$).
For RE at $N{=}16$, SJ ($\rho{=}0.281$), FK ($\rho{=}0.279$), and logprob ($\rho{=}0.287$) all plateau at comparable levels, suggesting that compound entity-then-relation errors impose a ceiling that additional samples cannot overcome.

\paragraph{Scale ablation: 4B vs.\ 8B.}
To isolate whether the correlation-selection gap is an artifact of model capacity, we replicate the core evaluation with Qwen3-4B, a model with half the parameters of our primary 8B model.

\begin{table}[t]
\centering
\small
\caption{Scale ablation: 4B vs.\ 8B on quality estimation ($\rho$) and selection (F1). All results use $N{=}8$, $T{=}1.0$. 8B values reproduced from Tables~\ref{tab:quality_estimation}--\ref{tab:selection}.}
\label{tab:scale-ablation}
\begin{tabular}{lcccccc}
\toprule
& \multicolumn{2}{c}{\textbf{CoNLL NER}} & \multicolumn{2}{c}{\textbf{SciERC NER}} & \multicolumn{2}{c}{\textbf{SciERC RE}} \\
\cmidrule(lr){2-3} \cmidrule(lr){4-5} \cmidrule(lr){6-7}
& 8B & 4B & 8B & 4B & 8B & 4B \\
\midrule
SJ $\rho$ & .436 & .422 & .360 & .422 & .250 & .267 \\
FK $\rho$ & .400 & .418 & .267 & .326 & .275 & .359 \\
\midrule
Greedy F1 & .912 & .923 & .609 & .662 & .374 & --- \\
Oracle F1 & ---  & .940 & .769 & .784 & --- & --- \\
\bottomrule
\end{tabular}
\end{table}

Table~\ref{tab:scale-ablation} reveals that quality estimation correlations at 4B are preserved or even strengthened: SciERC NER SJ rises from $\rho{=}0.360$ (8B) to $\rho{=}0.422$ (4B), a $+6.2$~pp increase, and FK improves from $\rho{=}0.267$ to $\rho{=}0.326$.
Yet the selection gap persists: on SciERC NER, all 4B signals perform \textit{worse} than greedy (SJ $-1.49$~pp, FK $-1.56$~pp, LP $-2.23$~pp), despite 13~pp of oracle headroom.
On CoNLL, 4B oracle headroom is only 1.70~pp, leaving all signals within $\leq$0.1~pp of greedy---marginal by any standard.
The gap is therefore independent of model capacity; task difficulty and per-sample variance remain the primary drivers.

\paragraph{Free-form decoding ablation.}
One concern is that constrained decoding (XGrammar) artificially compresses log-probability distributions, creating the tied-score phenomenon.
We test this by running SciERC NER inference without grammar constraints, relying solely on the fine-tuned model's learned output format.
The fine-tuned model achieves a 99.95\% valid JSON parse rate without constraints, confirming that grammar enforcement is unnecessary for well-tuned models.
Crucially, LP standard deviation is virtually unchanged (0.0206 $\to$ 0.0199, ratio 0.97), and the tied fraction drops by only 2~pp (59.7\% $\to$ 57.7\%).
NER selection gain decreases from $+0.59$~pp to $+0.18$~pp---both negligible (Table~\ref{tab:freeform-ablation} in the Appendix).
Score compression is therefore intrinsic to the model and task, not an artifact of constrained decoding.
This reinforces the universality of the correlation-selection gap: relaxing decoding constraints does not restore log-probability's discriminative power for selection.

\paragraph{Cross-model transfer.}
LLaMA-3.1-8B-Instruct replicates the structural-over-surface ordering on SciERC NER (SJ $\rho{=}0.325{\pm}0.021$, FK $\rho{=}0.256{\pm}0.024$, 3-seed mean).
On CoNLL, the cross-model gap is minimal (Table~\ref{tab:quality_estimation}): SJ achieves $\rho{=}0.431$ (LLaMA) vs.\ 0.436 (Qwen), $\Delta{=}0.005$.
On conditional RE (non-zero F1 instances), SJ recovers to $\rho{=}0.405$ while FK drops to $\rho{=}0.259$, confirming that structural advantage is robust across both models and evaluation conditions.

\paragraph{Evaluation circularity.}
Leave-one-out SJ (LOO-SJ), which excludes the target sample when computing confidence, yields $\Delta\rho{=}{-}0.014$ relative to standard SJ, with a maximum per-instance $\Delta\rho$ of $3.1{\times}10^{-4}$.
Evaluation circularity is negligible.

\paragraph{Conditional filtering reveals signal divergence.}
We define instances as \textit{tied} when the within-sample log-probability range falls below $\varepsilon{=}0.05$ nats; results are robust to $\varepsilon \in \{0.01, 0.05, 0.10\}$ (tied fractions of 45\%, 60\%, 72\% on NER, with consistent patterns across all thresholds).
Filtering to non-zero F1 instances (oracle $>$ 0) improves RE correlations (SJ $\rho{=}0.316$ from 0.250 at $N{=}8$) while NER remains stable ($\rho{\approx}0.38$).
At $N{=}16$, signal families diverge sharply on conditional RE: SJ maintains $\rho{=}0.308$, while logprob drops to $\rho{=}{-}0.026$---effectively zero---showing that logprob's apparent RE utility at $N{=}8$ does not generalize to larger sample sets.
For applications that can first filter out catastrophic failures, SJ's effective ranking approaches $\rho{\approx}0.4$ on both NER and RE, a regime where quality-aware aggregation can provide meaningful benefit.

\paragraph{Binary filtering vs.\ continuous ranking.}
The logprob paradox observed in Table~\ref{tab:quality_estimation}---highest AUROC (0.761) but lowest $\rho$ (0.307) on CoNLL with LLaMA---reflects a fundamental distinction in how signals encode quality information.
Log-probability provides coarse separation between reliable and unreliable instances (suited for accept/reject thresholds), while structural signals like SJ and Voting provide finer-grained ranking (suited for prioritization and quality-aware aggregation).
Figure~\ref{fig:reliability} visualizes this complementarity across the difficulty spectrum.
Practitioners should select signal type based on the downstream application: binary filtering favors logprob; ranking-dependent decisions favor structural consistency.

\begin{figure}[t]
\centering
\includegraphics[width=\columnwidth]{figures/fig_reliability_1024}
\caption{Reliability diagrams for confidence signals on SciERC NER. Structural signals (VC, SJ) show better calibration than log-probability, which exhibits overconfidence in the low-quality regime.}
\label{fig:reliability}
\end{figure}
